\section{Introduction}\label{sec:intro}

\subsection{The results}\label{sec:intro:results}

Let $\zeta(s)$ be the Riemann zeta function, and for an integer $N\ge1$ and a Dirichlet character
$\chi$ modulo $N$ let $L(s,\chi)$ be the meromorphic continuation of $\sum_{n\ge1}\chi(n)n^{-s}$.
We prove the following theorems.

\begin{theorem}[The Riemann Hypothesis]\label{thm:main-rh}
Let $s\in\C$ with $\zeta(s)=0$. If $s$ is not one of the trivial zeros $-2,-4,-6,\dots$, then
$\Re s=\frac12$.
\formal{riemannHypothesis}
\end{theorem}

\begin{theorem}[The Generalized Riemann Hypothesis for Dirichlet $L$-functions]\label{thm:main-grh}
Let $N\ge1$, let $\chi$ be a Dirichlet character modulo $N$, primitive or not, and let $s\in\C$ with
$\Re s>\frac12$. Then $L(s,\chi)\neq0$, unless $\chi$ is the trivial character and $s=1$.
\formal{LFunction_ne_zero_of_half_lt_re}
\end{theorem}

\begin{corollary}\label{cor:main-strip}
Let $\chi$ be a Dirichlet character modulo $N\ge1$ and $s\in\C$ with $0<\Re s<1$ and $L(s,\chi)=0$.
Then $\Re s=\frac12$. In particular, all non-trivial zeros of $L(s,\chi)$ for primitive $\chi$ lie on
the critical line.
\formal{LFunction_eq_zero_re} \formal{LFunction_eq_zero_re_of_isPrimitive}
\end{corollary}

The statements are those of the formal development, which uses Mathlib's definitions
\cite{Mathlib2020}. In particular, $L(s,\chi)$ is Mathlib's \texttt{DirichletCharacter.LFunction}:
it is entire for $\chi\neq\one$, and for the trivial character modulo $1$ it equals $\zeta(s)$.
Mathlib assigns to $\zeta$ the value $\frac12(\gamma-\log4\pi)$ at its pole, which is why $s\neq1$
appears as a hypothesis in the formal version of \cref{thm:main-rh}. The deduction of
\cref{thm:main-rh} and \cref{cor:main-strip} from \cref{thm:main-grh} is classical and is given in
\cref{sec:reduction}.

By the classical criteria of Robin \cite{Robin1984} and Lagarias \cite{Lagarias2002},
\cref{thm:main-rh} implies that $\sigma(n)<e^{\gamma}n\log\log n$ for all $n>5040$, and that
$\sigma(n)\le H_n+e^{H_n}\log H_n$ for all $n\ge1$, where $\sigma$ is the sum-of-divisors function and
$H_n$ the $n$-th harmonic number. Both inequalities have also been obtained formally, by combining
the formal proof with an independent formalization of these criteria \cite{Robin1984Lean}; see
\cite{Sutherland2026} and \cref{sec:verification:arith}.

\subsection{The formal proof}\label{sec:intro:formal}

The complete proof of \cref{thm:main-rh,thm:main-grh} and \cref{cor:main-strip} has been formalized
in the proof assistant Lean~4 \cite{deMouraUllrich2021} on top of Mathlib \cite{LeanSource2026}.
The formal development consists of $175\,187$ files with about $1.07$ million declarations. About
$720\,000$ of these declarations, with eleven million lines of source text, are used in the proof of
the main theorems. The formal statements of the main theorems refer only to Mathlib's definitions of
$\zeta$, of Dirichlet characters, of their $L$-functions and of primitivity.

The formal proof has been checked independently by A.~V.~Sutherland \cite{Sutherland2026}. All
modules compile without errors. The five main theorems depend only on the three standard axioms
\texttt{propext}, \texttt{Classical.choice} and \texttt{Quot.sound}. Every module has been replayed
through the Lean kernel. The exported proof is accepted by the independently written kernels
\emph{nanoda} and \emph{con-leche}, and \emph{lean4lean} accepts every compiled module. A sandboxed
rebuild reproduces the exported proof byte for byte. At the time of writing, two further
single-threaded kernel runs on the exported proof were still in progress. We refer to
\cref{sec:verification} for details.

\subsection{Strategy}\label{sec:intro:strategy}

The proof is modelled on Stepanov's elementary proof of the Riemann Hypothesis for curves over finite
fields \cite{Stepanov1969,Bombieri1974}; compare \cite{Weil1948,Deligne1974}. To bound the number of
points of a curve $C$ over $\F_q$, Stepanov constructs an auxiliary function that vanishes to high
order $m$ at every $\F_q$-rational point. Comparing the total order of vanishing with the degree gives
an upper bound $\#C(\F_q)\le q+O(\sqrt q)$, and the optimal choice of $m$ is of size $\sqrt q$. A
one-sided bound suffices, because it can be applied to all twists of $C$ at once.

Over the integers there is no evident notion of multiplicity. In the present argument the
multiplicity is supplied by chains
\[
\Uni(X_0)\longrightarrow\Uni(X_1)\longrightarrow\cdots\longrightarrow\Uni(X_m)
\]
of theta lifts between totally definite unitary groups attached to a CM field. The length of the
chain is comparable to the square root of the length of the prime sum being estimated
\cite{Announcement2026}. The argument has five steps.

\begin{enumerate}[label=(\arabic*)]
\item \emph{One-sided bounds suffice} (\cref{sec:reduction,sec:landau}). Theorem~B follows from the
statement that $L(s,\psi)\ne0$ for primitive $\psi$ and $\Re s>\frac12$. By a Landau-type theorem on
windows (\cref{thm:landau-window}), it is enough to prove an \emph{upper} bound
$G_\psi(y)\le y^{1/2+\eps}$ for the Riesz-weighted prime sums
\[
G(y)=\sum_{n\le y}\Lambda(n)\Big(\frac yn-1\Big)-y(\log y-1-\gamma),\qquad
g_\chi(y)=\sum_{n\le y}\Lambda(n)\chi(n)\Big(\frac yn-1\Big)+y\,\frac{L'}{L}(1,\chi),
\]
combined as in \cref{def:landau:Gm}. The bound is only needed for $y$ in windows $[Y,Y^\lambda]$
with $Y$ arbitrarily large.
\item \emph{Prime counts in Frobenius classes} (\cref{sec:hosts}). Fix a primitive character $\psi$.
At each scale $x$ we choose a \emph{host field}: an abelian number field $E$ with Galois group
$(\Z/2)^2\times A$, cut out by $\psi$, by the quadratic characters of two auxiliary primes $q$ and
$d$, and by $2$-power congruences. An exact orthogonality identity, the \emph{balance identity}
(\cref{prop:balance}), expresses weighted counts of the primes that are inert in $\Q(\sqrt{-d})$ and
lie in a given Frobenius class of $E$ in terms of the prime sums $g_\chi$ of the characters of $E$.
An upper bound with error $x^{1/2+O(\eps)}$ for these counts, averaged over a suitable set of
classes, implies the bound in (1) (\cref{thm:cheb-windows}).
\item \emph{The Stepanov-type bound} (\cref{sec:stepanov,sec:families}). The required upper bound
for the inert-prime counts comes from numerical records attached to lattices constructed from chains
of theta lifts of length $m\asymp\sqrt x$. Two families of inequalities control these records:
  \begin{itemize}
  \item \emph{local} inequalities: every inert prime $p$ contributes divisibility, measured by local
  lattice lengths at $p$, in proportion to the multiplicity $m$ (\cref{thm:linkG});
  \item \emph{archimedean} inequalities: slopes and covolumes, computed through the Siegel--Weil
  formula, mass formulas and the Chowla--Selberg formula, bound the total divisibility, as in the
  slope inequalities of Bost \cite{Bost1996,Bost2001}.
  \end{itemize}
Approximate orthogonality of the relevant theta series (Theorem~S, \cref{thm:theorem-S}), Green's
parametrization of the cuspidal characters of finite general linear groups \cite{Green1955}, and a
Galois-cohomological link complete the inputs. Combining these inequalities by linear arithmetic
yields the required upper bound (\cref{prop:gap-bound,thm:global-chain}).
\item \emph{Existence of hosts and Galois cohomology} (\cref{sec:top,sec:window}). Hosts with the
required properties exist at all scales in suitable windows. This uses a bound, the \emph{window
bound}, for the lengths of the Tate--Shafarevich groups $\Sha^1(G_{E,S_\ell},\Ql/\Zl(-j))$ of the host
fields (\cref{thm:vwin}). This bound is a weak form of Vandiver's conjecture, and it is proved using
Galois representations attached to polarized cuspidal representations (\cref{sec:galois}), base
change, and Kato-type Euler systems.
\item \emph{Covolumes and literature} (\cref{sec:covolume,sec:literature}). Arithmetic-geometric
covolume bounds and a large body of classical results enter the construction. All of them are proved
in the formal development.
\end{enumerate}

\subsection{Organisation and conventions}\label{sec:intro:conv}

\Cref{sec:reduction,sec:landau,sec:hosts} contain complete proofs. \Cref{sec:top} reduces the main
theorem to three propositions: Green's theorem, the window bound, and the non-emptiness of a
hypothesis package. \Cref{sec:stepanov} explains the Stepanov-type gap bound and how it follows from
a list of inequalities. \Cref{sec:families,sec:window,sec:galois,sec:covolume} describe how the
arithmetic inputs are established, and \cref{sec:literature} lists the published results that are
used. These sections form Part~I. Part~II proves the arithmetic inputs.

Most arithmetic inputs have formal proofs of hundreds of thousands of lines. Part~I states these results
precisely and gives outlines of their proofs; such results are marked with a superscript~\Fmark{}. Their
proofs are given in Part~II, converted from the formal development; \cref{sec:pf:intro} explains the
organisation and the conventions of Part~II, and \cref{tab:pf:map} lists the section of Part~II in
which each of these results is proved. Formal counterparts are cited by the name of the Lean
declaration in square brackets, as in \formal{UpperTheory.FinalTheorems.everyCharacter}.
\Cref{sec:dictionary} lists the formal counterparts of all numbered results of Part~I. When a formal
statement differs from the classical result its name suggests (for instance, when a weaker form
suffices), this is stated in a remark.
